\begin{proof}[Proof of \cref{obj:prop:rademacher-mixture-separation}]
\begin{enumerate}
\item Fix a row \(n\) and a realized design \(w=(\mathbf A_n,\mathbf Z_n)\).  Under the independent-sign prior, define the selected sign field
\[
\xi_i^{w}:=\xi_{i,\alpha_w(i)},
\qquad
\alpha_w(i):=\mathbf 1\{i\text{ is observed in a treated realized group under }w\}.
\]
Because the realized groups are pairwise disjoint, every observed unit has \(\alpha_w(i)\) equal to its realized treatment arm.  Hence the observation produced from the independent schedule is exactly the observation produced from the common-sign schedule with signs \((\xi_i^w)_{i\in\mathcal U_n}\).  For fixed \(w\), the vector \((\xi_i^w)_i\) has the same product Rademacher law as the common-sign vector.  Finite Fubini then gives, for every statistic \(\varphi\),
\[
\mathbb E_{\mathbb H_n^{\mathrm{same}}}\mathbb E_{\mathbf Z_n}\{\varphi(\mathcal O_n)\}
=
\mathbb E_{\mathbf Z_n}\mathbb E_{\mathbb H_n^{\mathrm{same}}}\{\varphi(\mathcal O_n)\}
=
\mathbb E_{\mathbf Z_n}\mathbb E_{\mathbb H_n^{\mathrm{ind}}}\{\varphi(\mathcal O_n)\}
=
\mathbb E_{\mathbb H_n^{\mathrm{ind}}}\mathbb E_{\mathbf Z_n}\{\varphi(\mathcal O_n)\}.
\]
% lean: CausalSmith.Experimentation.DenseGroupPartitionProjectionPhase.sameObservedExpectation_eq_independent

\item For any deterministic vector \(\boldsymbol x=(x_i)_{i\in\mathcal U_n}\), write
\[
\bar x_n:=\frac1n\sum_{i\in\mathcal U_n}x_i,
\qquad
\operatorname{PV}_n(\boldsymbol x):=\frac1{n-1}\sum_{i\in\mathcal U_n}(x_i-\bar x_n)^2,
\qquad
m_{\boldsymbol x}(Q):=\frac1M\sum_{i\in Q}x_i
\]
for \(Q\in\Omega_{n,M}\).  Sampling without replacement gives
\[
\operatorname{Var}_{Q\sim\mathrm{Unif}(\Omega_{n,M})}\{m_{\boldsymbol x}(Q)\}
=
\left(\frac1M-\frac1n\right)\operatorname{PV}_n(\boldsymbol x).
\]
If \(x_i\in\{-1,1\}\), then
\[
\operatorname{PV}_n(\boldsymbol x)
=
\frac n{n-1}\left(1-\bar x_n^2\right).
\]
Under the common-sign prior, \(\mathbb E\bar\xi_n=0\) and
\[
\operatorname{Var}(\bar\xi_n)=\frac1n,
\]
so \(\bar\xi_n\to 0\) in probability.  Since \cref{obj:ass:group-count-growth} implies \(n\to\infty\),
\[
\left(\frac1M-\frac1n\right)\frac n{n-1}\to\frac1M,
\]
and therefore, for each arm \(z\),
\[
V_{z,n}\stackrel{p}{\longrightarrow}\frac1M
\qquad
\text{under }\mathbb H_n^{\mathrm{same}}.
\]
% lean: CausalSmith.Experimentation.DenseGroupPartitionProjectionPhase.scheduleArray_same_armVar_tendstoInProb

\item The same calculation applies arm by arm under \(\mathbb H_n^{\mathrm{ind}}\).  For each fixed arm \(z\),
\[
\bar\xi_{z,n}:=\frac1n\sum_{i\in\mathcal U_n}\xi_{i,z},
\qquad
\mathbb E\bar\xi_{z,n}=0,
\qquad
\operatorname{Var}(\bar\xi_{z,n})=\frac1n,
\]
so \(\bar\xi_{z,n}\to0\) in probability and
\[
V_{z,n}\stackrel{p}{\longrightarrow}\frac1M
\qquad
\text{under }\mathbb H_n^{\mathrm{ind}}.
\]
% lean: CausalSmith.Experimentation.DenseGroupPartitionProjectionPhase.scheduleArray_independent_armVar_tendstoInProb

\item Under \(\mathbb H_n^{\mathrm{same}}\), the two arm tables coincide pointwise:
\[
h_{1,n}(Q)-h_{0,n}(Q)=0
\qquad
(Q\in\Omega_{n,M}).
\]
By linearity of the Johnson projection,
\[
E_{\tau,1,n}
=
\left\|\Pi_{n,1}(h_{1,n}-h_{0,n})\right\|_n^2
=
\|\Pi_{n,1}0\|_n^2
=
0.
\]
Thus \(E_{\tau,1,n}\stackrel{p}{\to}0\) under \(\mathbb H_n^{\mathrm{same}}\).
% lean: CausalSmith.Experimentation.DenseGroupPartitionProjectionPhase.scheduleArray_same_degreeOne_tendstoInProb

\item Under \(\mathbb H_n^{\mathrm{ind}}\), put
\[
d_i:=\xi_{i,1}-\xi_{i,0},
\qquad
\bar d_n:=\frac1n\sum_i d_i,
\qquad
\bar c_n:=\frac1n\sum_i \xi_{i,1}\xi_{i,0}.
\]
The arm-difference table is the slice sample mean
\[
h_{1,n}(Q)-h_{0,n}(Q)=m_{\boldsymbol d}(Q).
\]
This function is a linear combination of degree-one inclusion indicators; after subtracting its slice mean, it is orthogonal to constants and lies in \(\mathcal H_{n,1}\).  Hence the degree-one projection equals the centered sample mean, and
\[
E_{\tau,1,n}
=
\left(\frac1M-\frac1n\right)\operatorname{PV}_n(\boldsymbol d).
\]
Since
\[
\operatorname{PV}_n(\boldsymbol d)
=
\frac n{n-1}
\left(2-2\bar c_n-\bar d_n^2\right),
\]
while \(\bar c_n\to0\), \(\bar\xi_{1,n}\to0\), and \(\bar\xi_{0,n}\to0\) in probability, we get
\[
E_{\tau,1,n}\stackrel{p}{\longrightarrow}\frac2M
\qquad
\text{under }\mathbb H_n^{\mathrm{ind}}.
\]
% lean: CausalSmith.Experimentation.DenseGroupPartitionProjectionPhase.scheduleArray_independent_degreeOne_tendstoInProb

\item By \cref{obj:ass:stable-treatment-fraction},
\[
p_n\to p,\qquad 0<p<1,
\]
and hence
\[
\frac1{p_n}\to\frac1p,
\qquad
\frac1{1-p_n}\to\frac1{1-p}.
\]
Consequently, for either prior,
\[
\Theta_n:=
\frac{V_{1,n}}{p_n}+\frac{V_{0,n}}{1-p_n}
\stackrel{p}{\longrightarrow}
\frac1M\left(\frac1p+\frac1{1-p}\right)
=
\frac1{M p(1-p)}.
\]
% lean: CausalSmith.Experimentation.DenseGroupPartitionProjectionPhase.rademacher_mixture_separation

\item Since \(0<G_{1n}<G_n\), every row has \(G_n\ge2\), and therefore \(2M\le M G_n\le n\).  The degree-one Kneser eigenvalue in \cref{obj:thm:exact-kneser-identity} is
\[
\lambda_{n,1}=-\frac{M}{n-M}.
\]
Thus \(\lambda_{n,1}\to0\).  Writing \(N_n=M G_n\) for the grouped-unit count and using \cref{obj:ass:sampling-fraction},
\[
G_n\lambda_{n,1}
=
-\frac{M G_n}{n-M}
=
-\frac{N_n/n}{1-M/n}
\longrightarrow
-\rho .
\]
% lean: CausalSmith.Experimentation.DenseGroupPartitionProjectionPhase.kneserEigenvalue_one

\item The exact rowwise variance identity is used on the eventual rows where both treatment arms contain at least two realized groups.  These rows are eventual because \(p_n\to p\in(0,1)\) and \(G_n\to\infty\): eventually \(p_n>p/2\) and \(G_n>4/p\), forcing \(G_{1n}\ge2\), and eventually \(p_n<(1+p)/2\) and \(G_n>4/(1-p)\), forcing \(G_{0n}\ge2\).

On those rows, \cref{obj:thm:exact-pame-variance,obj:thm:exact-kneser-identity} give
\[
G_n\sigma_n^2
=
\Theta_n+G_n\lambda_{n,1}E_{\tau,1,n}
-\lambda_{n,1}\Theta_n
=
(1-\lambda_{n,1})\Theta_n
+
G_n\lambda_{n,1}E_{\tau,1,n}.
\]
Combining the convergence statements above yields
\[
G_n\sigma_n^2
\stackrel{p}{\longrightarrow}
\frac1{M p(1-p)}
\qquad
\text{under }\mathbb H_n^{\mathrm{same}},
\]
and
\[
G_n\sigma_n^2
\stackrel{p}{\longrightarrow}
\frac1{M p(1-p)}-\frac{2\rho}{M}
\qquad
\text{under }\mathbb H_n^{\mathrm{ind}}.
\]
% lean: CausalSmith.Experimentation.DenseGroupPartitionProjectionPhase.scaledSigmaSq_eq_of_additive_armTables

\item The separation is positive because \(M>0\) and \(\rho>0\):
\[
0<\frac{2\rho}{M}.
\]
Also \cref{obj:ass:sampling-fraction} gives \(\rho\le1\), and \(p(1-p)\le1/4\) follows from \((p-1/2)^2\ge0\).  Since \(p(1-p)>0\),
\[
\frac{1}{M p(1-p)}-\frac{2\rho}{M}
=
\frac1M\left(\frac1{p(1-p)}-2\rho\right)
\ge
\frac1M(4-2\rho)
\ge
\frac2M .
\]
% lean: CausalSmith.Experimentation.DenseGroupPartitionProjectionPhase.rademacher_mixture_separation

\item Finally, if \(X_n\to c\) in probability and \(t<c\), then
\[
\Pr\{X_n<t\}
\le
\Pr\{|X_n-c|>c-t\}
\to0,
\qquad
\Pr\{X_n\ge t\}\to1.
\]
Apply this with \(X_n=G_n\sigma_n^2\).  The same-sign limit exceeds the independent-sign limit by \(2\rho/M\), so every
\[
0<c_\sigma<\frac{1}{M p(1-p)}-\frac{2\rho}{M}
\]
lies below both limits.  Therefore
\[
\Pr_{\mathbb H_n^{\mathrm{same}}}\{G_n\sigma_n^2\ge c_\sigma\}\to1,
\qquad
\Pr_{\mathbb H_n^{\mathrm{ind}}}\{G_n\sigma_n^2\ge c_\sigma\}\to1.
\]
% lean: CausalSmith.Experimentation.DenseGroupPartitionProjectionPhase.probability_ge_tendsto_one_of_tendstoInProb
\end{enumerate}
\end{proof}
